\documentclass[11pt]{article}
\usepackage[T1]{fontenc}
\usepackage[margin=1in]{geometry}
\usepackage{amsmath,amssymb,amsthm,hyperref}
\newtheorem{theorem}{Theorem}
\newtheorem{lemma}[theorem]{Lemma}
\newtheorem{corollary}[theorem]{Corollary}
\theoremstyle{remark}
\newtheorem{remark}[theorem]{Remark}
\newcommand{\Q}{\mathbb Q}
\newcommand{\C}{\mathbb C}
\newcommand{\Z}{\mathbb Z}
\newcommand{\vp}{v_p}
\title{Irreducibility reductions for consecutive binomial moments}
\author{}
\date{}
\begin{document}
\maketitle
\begin{abstract}
Consecutive normalized moment polynomials of a univariate binomial never
divide one another in the direction of increasing degree, for all
nonnegative integral lower exponents, positive integral gaps, and positive
real Gamma shapes. Thus irreducibility of either polynomial suffices for
consecutive coprimality. For ordinary factorial moments we prove
factor-degree divisibility under an explicit no-carry condition. It yields
irreducibility at prime-power orders $p^r$ when $p$ exceeds the total support
degree, and a further coprimality criterion using adjacent prime-power parts. An independent Eisenstein argument gives
irreducibility over the rational function field of the Gamma shape and
consecutive coprimality at every transcendental shape. No all-order
irreducibility or coprimality claim is made at a fixed algebraic shape.
\end{abstract}

For integers $a\ge0$, $d\ge1$ and $n\ge1$, and a positive real number
$\beta$, put
\begin{equation}\label{eq:definition}
 P_{n,a,d,\beta}(t)
 =\sum_{k=0}^n\binom nk(an+\beta)_{dk}t^k,
 \qquad
 (x)_j=\prod_{i=0}^{j-1}(x+i),\quad (x)_0=1 .
\end{equation}
We abbreviate this polynomial by $P_n$ when the parameters are fixed.
It has constant term one and exact degree $n$.
For the functional $\mathcal L_\beta(X^j)=(\beta)_j$ and a genuine
binomial $F(X)=X^a(A+BX^d)$, $A,B\in\C^\times$, one has
\begin{equation}\label{eq:moment-normalization}
 \mathcal L_\beta(F^n)=A^n(\beta)_{an}P_n(B/A).
\end{equation}
For positive $\beta$, the prefactor is nonzero; $\beta=1$ gives the
ordinary factorial functional. Consecutive coprimality of $P_n,P_{n+1}$
therefore rules out two consecutive zero moments for every binomial
with the given support.

\section{A strict leading-coefficient separation}

\begin{theorem}[Consecutive nondivisibility]\label{thm:nondivisibility}
For all integers $a\ge0$, $d\ge1$, $n\ge1$ and all real $\beta>0$,
the polynomial $P_n$ does not divide $P_{n+1}$ in $\C[t]$.
More precisely, if $c_n=[t]P_{n+1}-[t]P_n$ and
$L_n=[t^{n+1}]P_{n+1}/[t^n]P_n$, then $L_n>c_n$.
\end{theorem}
\begin{proof}
Set
\[
 x=an+\beta>0,\qquad r=a+d,\qquad
 A_0=(x)_r,\qquad B_0=(x)_a(x)_d,\qquad f(s)=(x+s)_r .
\]
Direct comparison gives
\begin{align}
 c_n&=(n+1)(x+a)_d-n(x)_d,\label{eq:first-difference}\\
 L_n&=\frac{(x+dn)_r}{(x)_a}.\label{eq:leading-ratio}
\end{align}
For the second identity, cancel the common factors in
$(x+a)_{d(n+1)}/(x)_{dn}$, or use the identity
$(x)_{a+dn+d}=(x)_a(x+a)_{d(n+1)}=(x)_{dn}(x+dn)_r$.
Consequently,
\begin{equation}\label{eq:numerator}
 (x)_a(L_n-c_n)=f(dn)-(n+1)A_0+nB_0.
\end{equation}

The function $f$ is convex for $s\ge0$. Indeed, its factors are positive
affine functions, so its second derivative is a sum of nonnegative
products; when $r=1$, it is linear. Convexity at
$d=(1-1/n)0+(1/n)dn$ gives
\[
 f(dn)\ge A_0+n\bigl(f(d)-A_0\bigr).
\]
The inequality is valid also for $n=1$.
The elementary rising-factorial identities give
\[
 \frac{f(d)}{A_0}=\frac{(x+r)_d}{(x)_d}
       >\frac{(x+a)_d}{(x)_d}=\frac{A_0}{B_0},
\]
because $r=a+d>a$ and every factor is positive.
Thus
\begin{equation}\label{eq:strict-core}
 f(d)-2A_0+B_0
 >\frac{A_0^2}{B_0}-2A_0+B_0
 =\frac{(A_0-B_0)^2}{B_0}\ge0.
\end{equation}
Combining the last three displays shows
\[
 (x)_a(L_n-c_n)\ge n\bigl(f(d)-2A_0+B_0\bigr)>0.
\]
This proves the asserted strict inequality, including $a=0$ and $r=1$.

If $P_n$ divided $P_{n+1}$, their degrees and constant terms would make
the quotient $1+ct$. The coefficient of $t$ would force $c=c_n$,
while the leading coefficient would force $c=L_n$. This contradicts
$L_n>c_n$.
\end{proof}

\begin{corollary}[Irreducibility suffices]\label{cor:irreducible-pair}
Fix the parameters of Theorem~\ref{thm:nondivisibility}, and let
$K\subset\mathbb R$ be any subfield containing $\beta$.
If either $P_n$ or $P_{n+1}$ is irreducible in $K[t]$, then
$\gcd(P_n,P_{n+1})=1$ in $K[t]$ and in $\C[t]$.
In particular, this applies over $\Q$ whenever $\beta$ is rational.
The pair $P_1,P_2$ is coprime for every real $\beta>0$.
\end{corollary}
\begin{proof}
If $P_n$ is irreducible, a nonconstant gcd in $K[t]$ would make $P_n$
divide $P_{n+1}$, contrary to Theorem~\ref{thm:nondivisibility}.
If $P_{n+1}$ is irreducible, no polynomial of degree at most $n$ can
share a nonconstant factor with it over $K$.
A Bezout identity for a coprime pair over $K$ remains valid over $\C$.
Finally, $P_1$ has degree one; a shared complex root would make it
divide $P_2$, which the theorem excludes.
\end{proof}

\section{Ordinary moments and prime-power factor degrees}

The next theorem uses a prime-power part of the \emph{order}; the support
gap $d$ need not be a prime power.

\begin{theorem}[Factor degrees without a carry]\label{thm:factor-degrees}
Let $a\ge0$, $d\ge1$ be integers. Let $p$ be prime, $r\ge1$, $b\ge1$,
and suppose
\begin{equation}\label{eq:no-carry}
 (ab\bmod p)+db<p,
\end{equation}
where $ab\bmod p\in\{0,\ldots,p-1\}$. Put $P=p^r$ and $n=bP$.
Every root of the ordinary normalized moment polynomial $P_{n,a,d,1}$
in $\overline{\Q}_p$ has valuation
\begin{equation}\label{eq:root-valuation}
 -\frac{s}{P},\qquad
 s=d\,\frac{p^r-1}{p-1},\qquad \gcd(s,P)=1,
\end{equation}
where $\vp(p)=1$. Consequently every nonconstant factor over $\Q$
has degree divisible by $P$.
\end{theorem}
\begin{proof}
Write
\[
 P_{n,a,d,1}(t)=\sum_{k=0}^n c_kt^k,\qquad
 c_k=\binom nk(an+1)_{dk}.
\]
Condition~\eqref{eq:no-carry} implies $db<p$, and in particular
$b<p$ and $d<p$. It also says that $(ab,(a+d)b]$ contains no
multiple of $p$. Multiplication by $P=p^r$ shows that
$(an,(a+d)n]$ contains no multiple of $p^{r+1}$. Thus the shorter
interval $(an,an+dk]$ contains no multiple of $p^j$ for any $j>r$.
For $1\le j\le r$, the integer $an$ is divisible by $p^j$.
Legendre's formula therefore gives the exact identity
\begin{equation}\label{eq:shifted-valuation}
 \vp\bigl((an+1)_{dk}\bigr)
 =\sum_{j=1}^r\left(
     \left\lfloor\frac{an+dk}{p^j}\right\rfloor
       -\left\lfloor\frac{an}{p^j}\right\rfloor\right)
 =\sum_{j=1}^r\left\lfloor\frac{dk}{p^j}\right\rfloor
 =\vp((dk)!).
\end{equation}
The last equality uses $dk\le dbP<p^{r+1}$.
In particular, $\vp(c_0)=0$ and $\vp(c_n)=bs$.
Also $s=d(1+p+\cdots+p^{r-1})\not\equiv0\pmod p$.

For $0<k<n$, put $t=\vp(k)$, so $0\le t\le r$.
For every $1\le j<n$, one has $\vp(n-j)=\vp(j)$:
when $\vp(j)<r$ this follows from $\vp(n)=r$; otherwise
$j=\ell P$, $1\le\ell<b<p$, and both valuations equal $r$.
Consequently $\binom{n-1}{k-1}$ is a $p$-adic unit, and
\[
 \vp\binom nk=\vp\left(\frac nk\binom{n-1}{k-1}\right)=r-t.
\]
If $t<r$, the floors with $j\le t$ are exact and every remaining
floor is strictly greater than its argument minus one. Hence
\begin{equation}\label{eq:strict-chord}
 \vp(c_k)
 =r-t+\sum_{j=1}^r\left\lfloor\frac{dk}{p^j}\right\rfloor
 >dk\sum_{j=1}^r p^{-j}
 =\frac{sk}{P}.
\end{equation}
If $t=r$, all those arguments are integers, so
$\vp(c_k)=sk/P$. Thus every coefficient lies on or above the line
joining the two endpoint valuations.

If a root $\xi$ had $\vp(\xi)>-s/P$, the constant term of
$P_n(\xi)$ would have uniquely least valuation.
If $\vp(\xi)<-s/P$, the leading term would have uniquely least
valuation. The coefficient bounds just proved exclude cancellation
in either case. This proves \eqref{eq:root-valuation}, without needing
an additional Newton-polygon theorem.

For a rational factor of degree $h$, the quotient of its nonzero
constant coefficient by its leading coefficient has valuation
$-hs/P$, the valuation of the product of its $h$ roots. This is an
integer because the quotient is rational. Since $\gcd(s,P)=1$,
one has $P\mid h$.
\end{proof}

\begin{corollary}[Prime-power-order irreducibility]\label{thm:prime-power-order}
If $b=1$ in Theorem~\ref{thm:factor-degrees}, then
$P_{p^r,a,d,1}$ is irreducible over $\Q$.
Its sufficient no-carry condition is
\[
 (a\bmod p)+d<p.
\]
In particular, this holds whenever $p>a+d$.
\end{corollary}
\begin{proof}
Every nonconstant rational factor has degree divisible by $p^r$,
which is the degree of the whole polynomial.
\end{proof}

\begin{corollary}[Ordinary pairs from one irreducible order]\label{cor:ordinary-orders}
For fixed integers $a\ge0$, $d\ge1$, let $N\ge2$.
If either $N$ or $N-1$ is $p^r$ with $r\ge1$ and
$(a\bmod p)+d<p$, then
\[
 \gcd(P_{N-1,a,d,1},P_{N,a,d,1})=1.
\]
This holds in particular if the indicated base prime satisfies $p>a+d$.
The same conclusion holds for $N=2$ without any prime hypothesis.
No binomial $X^a(A+BX^d)$, $AB\ne0$, can have zero ordinary
moments at both orders in these cases.
\end{corollary}
\begin{proof}
Apply Corollary~\ref{thm:prime-power-order} to the indicated order and
Corollary~\ref{cor:irreducible-pair}. The case $N=2$ uses the last
assertion of the latter corollary. Equation~\eqref{eq:moment-normalization}
gives the moment conclusion.
\end{proof}

\begin{corollary}[Adjacent prime-power parts]\label{cor:adjacent-parts}
Fix $a\ge0$, $d\ge1$. Suppose that for primes $p,q$, positive integers
$r,s,b,c$, and some $n\ge1$,
\[
 n=bp^r,\qquad n+1=cq^s,\qquad
 (ab\bmod p)+db<p,\qquad (ac\bmod q)+dc<q.
\]
Then $\gcd(P_{n,a,d,1},P_{n+1,a,d,1})=1$ over $\C$.
The two orders need not themselves be prime powers.
\end{corollary}
\begin{proof}
Put $P=p^r$, $Q=q^s$. The primes are distinct, since they divide
consecutive integers. By Theorem~\ref{thm:factor-degrees}, a common
rational factor of degree $h$ would satisfy $P\mid h$ and $Q\mid h$.
The no-carry conditions give $P\ge p>db$ and $Q\ge q>dc$, hence
\[
 P^2>dn,\qquad Q^2>d(n+1),\qquad
 PQ>d\sqrt{n(n+1)}>n .
\]
Thus $PQ\mid h$ is impossible for $0<h\le n$.
A Bezout identity transfers coprimality to $\C$.
\end{proof}

\begin{remark}\label{rem:prime-power-scope}
For the ordinary support $\{2,5\}$, take
$n=145=5\cdot29$ and $n+1=146=2\cdot73$.
The no-carry sums are $10+15=25<29$ and $4+6=10<73$.
Corollary~\ref{cor:adjacent-parts} therefore applies although neither
order is a prime power.

For $a=0$, compare the valuation argument with Lemma~2.3 of
Filaseta and Rocks~\cite{FR}. The proof above includes the shift $an$
and proves the coefficient inequalities directly.
Irreducibility and consecutive coprimality at every order remain open
for a general fixed algebraic shape.
\end{remark}

\section{Generic and transcendental Gamma shapes}

Here the symbol $b$ is an indeterminate; it is not a fixed positive
real parameter.

\begin{theorem}[Generic-shape Eisenstein argument]\label{thm:generic-shape}
For every pair of integers $a\ge0$, $d\ge1$ and every $n\ge1$, the
polynomial
\[
 P_{n,a,d,b}(t)\in\Q[b][t]
\]
is irreducible over $\Q(b)$. The same holds for its monic reciprocal
\[
 R_{n,a,d,b}(u)
 =u^nP_{n,a,d,b}(u^{-1})
 =\sum_{k=0}^n\binom nk(an+b)_{dk}u^{n-k}.
\]
\end{theorem}
\begin{proof}
Use the prime element $\pi=b+an$ in $\Q[b]$.
Every nonleading coefficient of $R_n$ is divisible by $\pi$.
Its constant coefficient is
\[
 (\pi)_{dn}=\pi\prod_{j=1}^{dn-1}(\pi+j),
\]
which is not divisible by $\pi^2$; its quotient by $\pi$ has value
$(dn-1)!\ne0$ at $\pi=0$.

This is the Eisenstein argument over $\Q[b]$.
Explicitly, a factorization over $\Q(b)$ would, by Gauss' lemma,
give monic positive-degree factors in $\Q[b][u]$.
Their reductions modulo $\pi$ would be powers of $u$, since their
product reduces to $u^n$. Both constant coefficients would then be
divisible by $\pi$, forcing the constant coefficient of $R_n$ to be
divisible by $\pi^2$, a contradiction.
The nonzero constant and leading coefficients make reciprocal
transformation preserve reducibility over $\Q(b)$.
\end{proof}

\begin{corollary}[Transcendental shapes]\label{cor:transcendental-shape}
Let $\beta\in\C$ be transcendental over $\Q$.
For fixed integers $a\ge0$, $d\ge1$, every polynomial
$P_{n,a,d,\beta}$ is irreducible over $\Q(\beta)$, and
\[
 \gcd(P_{n,a,d,\beta},P_{n+1,a,d,\beta})=1
 \qquad(n\ge1)
\]
over $\C$.
For every positive real transcendental Gamma shape, this gives
coprimality at every consecutive pair of orders.
\end{corollary}
\begin{proof}
Evaluation $b\mapsto\beta$ is a field isomorphism
$\Q(b)\simeq\Q(\beta)$, so it preserves the irreducibility in
Theorem~\ref{thm:generic-shape}.
The consecutive polynomials are irreducible over the same field and
have different degrees, hence are coprime there.
Their Bezout identity remains valid over $\C$.
This conclusion uses neither specialization at an algebraic shape
nor an assertion of irreducibility over $\C$.
\end{proof}

\begin{remark}\label{rem:algebraic-specialization}
Theorem~\ref{thm:generic-shape} does not give irreducibility at a
prescribed rational or algebraic value of $b$. This distinction is
necessary: for $a=0$, $d=1$, $n=3$, $\beta=1/4$ one has
\[
 P_{3,0,1,1/4}(t)=\frac{(3t+4)(15t^2+16)}{64}.
\]
Thus it is reducible and has the rational root $-4/3$.
At that root the adjacent polynomials have values
$P_{2,0,1,1/4}(-4/3)=8/9$ and
$P_{4,0,1,1/4}(-4/3)=32/9$, so this factorization is not a
counterexample to consecutive coprimality.
In particular, generic irreducibility does not resolve the remaining
ordinary case $\beta=1$.
\end{remark}

\section{The support \texorpdfstring{$\{2,5\}$}{2,5}}

\begin{corollary}[Explicit quintic coefficient separation]\label{cor:quintic}
For $a=2$, $d=3$, $\beta=1$, one has
\[
 P_n(t)=\sum_{k=0}^n\binom nk\frac{(2n+3k)!}{(2n)!}t^k
\]
and the exact identity
\[
 L_n-c_n
 =\frac{3n(3n+2)(9n+7)(37n+17)}{2(2n+1)}>0 .
\]
Thus irreducibility of either consecutive polynomial suffices for
coprimality; this applies in particular when the corresponding
order is $p^r$ with $p>5$.
\end{corollary}
\begin{proof}
Here
\[
 c_n=4(n+1)(2n+3)(4n+5),\qquad
 L_n=\frac{5\prod_{j=1}^4(5n+j)}{2(2n+1)}.
\]
Subtracting and clearing the denominator gives
\[
 5\prod_{j=1}^4(5n+j)
 -8(2n+1)(n+1)(2n+3)(4n+5)
 =3n(3n+2)(9n+7)(37n+17).
\]
The remaining assertions follow from the preceding results.
For $n=1$ explicitly, $P_1=1+60t$,
$P_2=1+420t+151200t^2$, and $P_2(-1/60)=36$.
\end{proof}

The general nondivisibility proof is independent of irreducibility
theorems for ordinary Bessel polynomials. The latter are
$ {}_2F_0$ polynomials, whereas the displayed quintic family is
$ {}_4F_0$ with three varying positive upper parameters in addition
to $-n$. No Bessel irreducibility theorem is imported here.

\paragraph{Exact replay.}
The accompanying checker verifies the displayed algebraic identities,
the defining coefficient and prime-power valuation formulas on explicit
bounded inputs, and the generic Eisenstein coefficient conditions.
These are regression controls on the written proofs.
They do not constitute an all-order factorization search, an independent
replay, or formal verification.

\begin{thebibliography}{9}
\bibitem{FR}
M.~Filaseta and B.~Rocks,
\emph{On the irreducibility of a polynomial associated with the Strong
Factorial Conjecture},
Colloq. Math. \textbf{145} (2016), 307--314.
\href{https://doi.org/10.4064/cm6761-3-2016}{doi:10.4064/cm6761-3-2016}.
\end{thebibliography}
\end{document}
